% =====================================================================
% Глава 1. Аналитический обзор литературы.
% Объем 6--7 страниц.
% Структура по правкам друга-ревьюера
%   1.1 Мульти-агентные системы больших языковых моделей
%   1.2 Коммуникационные топологии агентов
%   1.3 Механизмы адаптации топологий
%   1.4 Безопасность и устойчивость топологий
%   1.5 Человек в контуре управления мульти-агентными системами
%   1.6 Сравнительная аналитическая таблица методов
%   1.7 Что отсутствует в существующих работах
%   1.8 Выводы по главе
% =====================================================================

\section{Аналитический обзор литературы}
\label{sec:review}

Настоящая глава представляет аналитический обзор работ, на стыке
которых формируется предметная область исследования. Обзор
охватывает пять направлений --- архитектуру мульти-агентных систем
на основе больших языковых моделей, коммуникационные топологии
агентов, механизмы их адаптации, безопасность и устойчивость
топологий, а также интеграцию человека в контур управления
мульти-агентными системами. Глава завершается сравнительной
таблицей основных методов и определением исследовательской лакуны,
заполнение которой составляет содержательный вклад настоящей
работы.

Методология обзора. В рассматриваемую выборку включены работы
с 2022 по 2026 годы, отобранные по ключевым словам \emph{multi-agent
LLM}, \emph{agent topology}, \emph{LLM debate}, \emph{human-in-the-
loop multi-agent} в базах публикаций arXiv, ACL Anthology, ICLR,
NeurIPS, ICML. Из первоначального списка из более чем шестидесяти
работ оставлены те, которые либо предлагают новую архитектуру
взаимодействия агентов, либо систематически анализируют свойства
существующих архитектур. Работы, посвященные исключительно
обучению одиночных моделей или классическому планированию в
дискретных средах, в обзор не включены.

\subsection{Мульти-агентные системы больших языковых моделей}
\label{sec:rev:mas}

Идея использования нескольких автономных программных агентов для
совместного решения задач восходит к классическим работам по
распределенному искусственному интеллекту. С появлением больших
языковых моделей (от англ. Large Language Models, LLM)
концепция мульти-агентных систем (от англ. Multi-Agent System,
MAS) переживает второе рождение --- роль агента теперь играет
языковая модель с заданной системной подсказкой и доступным
набором инструментов. Обзорная работа~\cite{wang2024} систематизирует
быстрорастущую совокупность исследований и выделяет несколько
устойчивых архитектурных образцов.

Первой широко известной системой подобного класса является
фреймворк AutoGen~\cite{wu2023}, в котором несколько агентов
обмениваются сообщениями в гибком формате многораундового диалога
и могут вызывать внешние инструменты, в том числе подключать
человека. AutoGen не фиксирует топологию заранее --- структура
взаимодействия определяется императивным кодом конкретного
сценария. Это позволяет реализовать любую конфигурацию, но
лишает фреймворк возможности сравнивать топологии между собой
в рамках единой архитектуры.

Иной образец продемонстрирован в системе ChatDev~\cite{qian2024},
где процесс разработки программного обеспечения моделируется
последовательностью ролевых агентов --- генерального директора,
программиста, тестировщика --- работающих по принципу
водопадной модели. Метод MetaGPT~\cite{hong2024} развивает
идею за счет введения стандартных операционных процедур, что
позволяет снизить число ошибок согласования между ролями.
Эти работы демонстрируют, что разделение труда между специализированными
агентами повышает качество решения сложных задач, однако
топология взаимодействия в них жестко зашита в архитектуру.

Параллельно развивается направление, в котором отдельный агент
итеративно улучшает собственный ответ. Метод Self-Refine
\cite{selfrefine2024} реализует цикл генерация --- критика ---
переработка, а метод Reflexion~\cite{reflexion2023} расширяет
цикл вербальным подкреплением --- модель формирует словесную
обратную связь и хранит ее в эпизодической памяти. Несмотря
на привлекательную простоту, эти методы по существу остаются
одноагентными и не используют разнообразие точек зрения,
доступное мульти-агентным конфигурациям.

\subsection{Коммуникационные топологии агентов}
\label{sec:rev:topology}

Под коммуникационной топологией понимается граф, описывающий,
какие агенты могут обмениваться сообщениями и в каком порядке.
В существующих фреймворках топология выбирается архитектурно
и не пересматривается во время выполнения задачи. Тем не менее
ряд недавних работ ставит вопрос об оптимальной структуре
такого графа. Используемая в настоящей работе таксономия графовых
структур включает пять базовых типов --- звезда (центральный
координатор и периферийные работники), цепочка (линейный конвейер),
полносвязная (каждый агент общается с каждым), дебатная (два
оппонента и арбитр) и иерархическая (двух- или многоуровневое
разделение труда). Указанные пять типов и составляют пространство
сравнения, исследуемое в настоящей работе. Схематическое
изображение пяти базовых топологий приведено на
рисунке~\ref{fig:rev:topologies}.

\begin{figure}[!htbp]
\centering
\begin{tikzpicture}[
  font=\scriptsize,
  node distance=0.35cm,
  every node/.style={align=center},
  ag/.style={circle, draw=black, thick, minimum size=4.5mm,
             inner sep=0.5pt, fill=blue!10},
  hub/.style={circle, draw=black, thick, minimum size=5mm,
              inner sep=0.5pt, fill=orange!25},
  jud/.style={circle, draw=black, thick, minimum size=5mm,
              inner sep=0.5pt, fill=green!18},
  ed/.style={-{Latex[length=1mm]}, thick},
  bid/.style={{Latex[length=1mm]}-{Latex[length=1mm]}, thick},
]

% --- Star ---
\begin{scope}[xshift=0cm]
  \node[hub] (s_c) at (0,0) {C};
  \node[ag] (s_1) at (90:0.95)  {1};
  \node[ag] (s_2) at (162:0.95) {2};
  \node[ag] (s_3) at (234:0.95) {3};
  \node[ag] (s_4) at (306:0.95) {4};
  \node[ag] (s_5) at (18:0.95)  {5};
  \foreach \i in {s_1,s_2,s_3,s_4,s_5}{ \draw[bid] (s_c) -- (\i); }
  \node[below=1.2cm of s_c, font=\footnotesize\bfseries] {Star};
\end{scope}

% --- Chain ---
\begin{scope}[xshift=3.0cm]
  \node[ag] (c_1) at (-1.0,0) {1};
  \node[ag] (c_2) at (-0.0,0) {2};
  \node[ag] (c_3) at (1.0,0)  {3};
  \draw[ed] (c_1) -- (c_2);
  \draw[ed] (c_2) -- (c_3);
  \node[below=1.1cm of c_2, font=\footnotesize\bfseries] {Chain};
\end{scope}

% --- Mesh ---
\begin{scope}[xshift=6.5cm]
  \node[ag] (m_1) at (90:0.85)  {1};
  \node[ag] (m_2) at (18:0.85)  {2};
  \node[ag] (m_3) at (306:0.85) {3};
  \node[ag] (m_4) at (234:0.85) {4};
  \node[ag] (m_5) at (162:0.85) {5};
  \foreach \i/\j in {m_1/m_2, m_1/m_3, m_1/m_4, m_1/m_5,
                     m_2/m_3, m_2/m_4, m_2/m_5,
                     m_3/m_4, m_3/m_5, m_4/m_5} {
    \draw[thick, gray!70] (\i) -- (\j);
  }
  \node[below=1.1cm of m_3, font=\footnotesize\bfseries] {Mesh};
\end{scope}

% --- Debate (вторая строка) ---
\begin{scope}[xshift=1.5cm, yshift=-3.4cm]
  \node[ag] (d_p) at (-0.8, 0.4) {$+$};
  \node[ag] (d_n) at (0.8, 0.4)  {$-$};
  \node[jud] (d_j) at (0, -0.7)  {J};
  \draw[bid] (d_p) -- (d_n);
  \draw[ed] (d_p) -- (d_j);
  \draw[ed] (d_n) -- (d_j);
  \node[below=0.4cm of d_j, font=\footnotesize\bfseries] {Debate};
\end{scope}

% --- Hierarchical ---
\begin{scope}[xshift=5.5cm, yshift=-3.4cm]
  \node[hub] (h_c) at (0, 0.6) {C};
  \node[ag] (h_a1) at (-1.0, -0.2) {a1};
  \node[ag] (h_a2) at (-0.3, -0.2) {a2};
  \node[ag] (h_b1) at (0.3, -0.2)  {b1};
  \node[ag] (h_b2) at (1.0, -0.2)  {b2};
  \draw[ed] (h_c) -- (h_a1);
  \draw[ed] (h_c) -- (h_a2);
  \draw[ed] (h_c) -- (h_b1);
  \draw[ed] (h_c) -- (h_b2);
  \draw[bid, gray!70] (h_a1) -- (h_a2);
  \draw[bid, gray!70] (h_b1) -- (h_b2);
  \node[below=0.85cm of h_c, font=\footnotesize\bfseries]
       {Hierarchical};
\end{scope}

\end{tikzpicture}
\caption{Схематическое изображение пяти базовых коммуникационных
топологий, рассматриваемых в работе. Узлы --- агенты; C ---
координатор, J --- арбитр, $+$ и $-$ --- оппоненты в дебатах,
a1, a2, b1, b2 --- участники подкоманд. Сплошные направленные
ребра обозначают передачу управления, серые двунаправленные ---
обмен сообщениями в общей шине.}
\label{fig:rev:topologies}
\end{figure}

Работа GPTSwarm~\cite{zhuge2024} рассматривает мульти-агентную
систему как обучаемый граф вычислений и оптимизирует его
параметры на корпусе задач. Авторы показывают, что
оптимизированные графы превосходят случайные конфигурации,
но в качестве пространства поиска используют только статичные
структуры. Метод G-Designer~\cite{zhang2024} применяет
графовые нейронные сети (от англ. Graph Neural Network ---
GNN) для автоматического подбора топологии под конкретную
задачу. Несмотря на адаптацию к типу задачи, выбор топологии
производится однократно и не пересматривается в ходе
решения.

В работе~\cite{qian2024scaling} систематически исследуются
закономерности масштабирования мульти-агентных систем при
увеличении числа агентов. Авторы обнаруживают, что топологии
типа малого мира (small-world) демонстрируют повышенную
робастность к ошибкам отдельных агентов, однако сравнение
ведется между статичными конфигурациями и не учитывает
динамику решения задачи. Метод DyLAN~\cite{liu2023}
предлагает динамический выбор состава команды агентов на
каждом шаге, но варьирует только участников --- сама
структура связей между ними остается фиксированной. Метод
Mixture-of-Agents~\cite{moa2024} организует агентов
в многослойную ансамблевую конфигурацию с агрегацией ответов,
которая также не меняется во время решения.

Промежуточные между однозвенными и мульти-агентными решения
представлены работами Tree-of-Thoughts~\cite{tot2023} и
CAMEL~\cite{camel2023}. Первая разворачивает дерево вариантов
рассуждения с поиском в ширину, вторая моделирует ролевой
диалог двух агентов. В обоих случаях структура взаимодействия
определяется заранее и фиксируется на протяжении всей задачи.
Ближе к проблеме адаптации структуры подходят работы по дебатам
между языковыми моделями~\cite{du2023}, где несколько моделей
обмениваются аргументами для повышения качества вывода, однако
число оппонентов и протокол обмена в них зафиксированы заранее.

Общая характеристика рассмотренных работ заключается в том,
что выбор топологии в них производится либо архитектурно
(на этапе проектирования системы), либо заранее (на этапе
получения новой задачи). Возможность переключения топологии
в процессе решения, в зависимости от фазы и наблюдаемых
сигналов, в этих работах не рассматривается.

\subsection{Механизмы адаптации топологий}
\label{sec:rev:adaptation}

Адаптация мульти-агентных систем во время выполнения исследуется
существенно реже. В проанализированной выборке преобладают два
направления адаптации, ни одно из которых не покрывает
интересующий нас случай динамического переключения топологии между
фазами решения.

Первое направление --- адаптация состава команды. Уже упоминавшийся
метод DyLAN~\cite{liu2023} выбирает на каждом шаге подмножество
агентов, наиболее подходящее для текущего контекста. Аналогично
работа AgentVerse~\cite{agentverse2023} использует механизм
найма агентов, при котором система формирует команду под задачу
из заранее заданного пула. В обоих случаях изменяется не структура
взаимодействия, а множество участников --- сама схема обмена
сообщениями остается фиксированной.

Второе направление --- оптимизация конфигурации перед запуском
задачи. Системы GPTSwarm~\cite{zhuge2024} и
G-Designer~\cite{zhang2024} принимают решение о структуре графа
агентов до начала вычислений и не пересматривают его в дальнейшем.
Такие методы можно рассматривать как однократную офлайн-адаптацию,
которая полезна для классов однородных задач, но недостаточна
для задач со сложной фазовой структурой.

Сложные задачи неоднородны по своей внутренней структуре. Фаза
планирования может требовать дебатного формата для генерации
альтернатив, фаза исполнения --- иерархического разделения
труда, фаза проверки --- независимой рецензии. Отсюда следует
естественная гипотеза --- статическая топология, оптимальная
для усредненной картины, может быть субоптимальной для каждой
отдельной фазы. Эта гипотеза мотивирует разработку
механизма переключения топологии в процессе выполнения задачи,
который и является одним из центральных предметов настоящего
исследования.

\subsection{Безопасность и устойчивость топологий}
\label{sec:rev:safety}

Отдельное направление работ посвящено анализу уязвимостей
мульти-агентных систем и зависимости этих уязвимостей от
топологии. Работа NetSafe~\cite{yu2024} систематически исследует
устойчивость различных конфигураций к воздействию скомпрометированного
агента. Авторы устанавливают, что плотные топологии (полносвязная,
звезда) более подвержены каскадному распространению ошибок, тогда
как разреженные структуры (цепочка) ограничивают радиус
поражения. Работа TOMA~\cite{liang2024} развивает анализ на
случай многошаговых атак, в которых злоумышленник может
последовательно воздействовать на несколько агентов с учетом
их связности.

Эти результаты важны для нашего исследования по двум причинам.
Во-первых, они показывают, что выбор топологии не сводится
к чистому критерию качества решения --- надежность системы
к ошибкам отдельных компонентов также является функцией
структуры графа. Во-вторых, они мотивируют включение в систему
защитных правил переключения топологий, которые предотвращают
патологические режимы взаимодействия (в частности, частое
осцилляционное переключение между топологиями без существенного
прогресса в решении). Реализация таких правил в настоящей работе
обсуждается в главе~\ref{sec:architecture}.

\subsection{Человек в контуре управления}
\label{sec:rev:hitl}

Подключение человека к работе автоматизированных систем (от англ.
Human-in-the-Loop, HITL) --- давно установленная практика. Обзорная
работа~\cite{mosqueira2023} систематизирует методы интеграции
человека в контур машинного обучения, выделяя такие функции
человека как разметка данных, исправление предсказаний, активный
отбор примеров. Классическая работа Хорвица~\cite{horvitz1999}
ввела принципы смешанной инициативы, при которой автоматизированная
система и человек поочередно ведут процесс решения задачи в
зависимости от уверенности каждой стороны. Эти принципы остаются
актуальными для проектирования современных мульти-агентных
систем.

Отдельную линию исследований представляют методы обучения
языковых моделей с обратной связью от человека. Работа
Кристиано и соавторов~\cite{christiano2017} предложила обучение
с подкреплением от человеческих предпочтений, что позднее легло
в основу метода InstructGPT~\cite{ouyang2022} и стало стандартной
техникой выравнивания моделей. В этих работах человек выступает
в роли поставщика сигнала вознаграждения. Применительно к
мульти-агентным системам вопрос ставится принципиально иначе ---
человек участвует не только в обучении, но и в каждом отдельном
вызове системы, что требует формализации его роли в графе
обмена сообщениями.

Обзор~\cite{wang2024} перечисляет несколько типичных паттернов
включения человека --- координатор, рецензент, судья при споре
агентов, --- но не предлагает экспериментального сравнения их
эффективности. Фреймворк AutoGen~\cite{wu2023} допускает
подключение человека как агента в диалоге, однако роль человека
задается архитектурно и не варьируется в ходе экспериментов.
Отдельные работы исследуют человека как судью в задачах оценки
качества генерации с использованием шкалы NASA-TLX (от англ.
NASA Task Load Index, индекс рабочей нагрузки NASA)~\cite{nasatlx1988},
но эта функция выделена в самостоятельную ветвь литературы и
не пересекается с проектированием рабочего графа агентов.

Современные оркестраторы мульти-агентных систем, такие как
Magentic-One~\cite{magenticone2024}, реализуют разделение между
ведущим агентом-координатором и подчиненными работниками, при
этом человек явно подключен как опциональный участник в роли
постановщика задачи. Однако сравнения эффективности разных ролей
человека в этой работе также не проводится.

В проанализированной выборке работ за 2022--2026 годы
систематических исследований того, какая роль человека и в какой
фазе наиболее эффективна для разных типов задач, не обнаружено.
Эта лакуна составляет второй центральный предмет настоящего
исследования.

\subsection{Сравнительная аналитическая таблица методов}
\label{sec:rev:comparison}

Рассмотренные методы естественно сопоставить по нескольким
критериям, релевантным для исследовательских вопросов работы.
Соответствующее сравнение приведено в таблице~\ref{tab:rev:methods}.
Знак \enquote{$+$} означает наличие свойства, \enquote{$-$} ---
его отсутствие, \enquote{$\pm$} --- частичную реализацию.

\begin{table}[!htbp]
\footnotesize
\caption{Сравнение существующих методов по ключевым критериям}
\label{tab:rev:methods}
\footnotesize
\begin{tabularx}{\textwidth}{l c c c c c c c}
\toprule
\textbf{Метод} & \textbf{Год} & \textbf{MAS} & \textbf{Сравн.}
& \textbf{Адапт.} & \textbf{Адапт.} & \textbf{HITL} & \textbf{Безопасн.} \\
& & & \textbf{топ.} & \textbf{состава} & \textbf{топ. runtime} & & \\
\midrule
ReAct~\cite{react2022}             & 2022 & $-$ & $-$ & $-$ & $-$ & $-$ & $-$ \\
AgentVerse~\cite{agentverse2023}   & 2023 & $+$ & $-$ & $+$ & $-$ & $-$ & $-$ \\
LLM debate~\cite{du2023}           & 2023 & $+$ & $-$ & $-$ & $-$ & $-$ & $-$ \\
CAMEL~\cite{camel2023}             & 2023 & $+$ & $-$ & $-$ & $-$ & $-$ & $-$ \\
DyLAN~\cite{liu2023}               & 2023 & $+$ & $-$ & $+$ & $-$ & $-$ & $-$ \\
Reflexion~\cite{reflexion2023}     & 2023 & $-$ & $-$ & $-$ & $-$ & $-$ & $-$ \\
AutoGen~\cite{wu2023}              & 2023 & $+$ & $-$ & $-$ & $-$ & $+$ & $-$ \\
Tree-of-Thoughts~\cite{tot2023}    & 2023 & $-$ & $-$ & $-$ & $-$ & $-$ & $-$ \\
Magentic-One~\cite{magenticone2024}& 2024 & $+$ & $-$ & $-$ & $-$ & $\pm$ & $-$ \\
MetaGPT~\cite{hong2024}            & 2024 & $+$ & $-$ & $-$ & $-$ & $-$ & $-$ \\
TOMA~\cite{liang2024}              & 2024 & $+$ & $\pm$ & $-$ & $-$ & $-$ & $+$ \\
Self-Refine~\cite{selfrefine2024}  & 2024 & $-$ & $-$ & $-$ & $-$ & $-$ & $-$ \\
ChatDev~\cite{qian2024}            & 2024 & $+$ & $-$ & $-$ & $-$ & $-$ & $-$ \\
Scaling MAS~\cite{qian2024scaling} & 2024 & $+$ & $\pm$ & $-$ & $-$ & $-$ & $\pm$ \\
MoA~\cite{moa2024}                 & 2024 & $+$ & $-$ & $-$ & $-$ & $-$ & $-$ \\
NetSafe~\cite{yu2024}              & 2024 & $+$ & $+$ & $-$ & $-$ & $-$ & $+$ \\
G-Designer~\cite{zhang2024}        & 2024 & $+$ & $+$ & $-$ & $-$ & $-$ & $-$ \\
GPTSwarm~\cite{zhuge2024}          & 2024 & $+$ & $\pm$ & $-$ & $-$ & $-$ & $-$ \\
\midrule
\textbf{Настоящая работа}          & 2026 & $+$ & $+$ & $-$ & $+$ & $+$ & $\pm$ \\
\bottomrule
\end{tabularx}
\\[0.3cm]
\footnotesize
Обозначения колонок ---
\textbf{MAS} мульти-агентная архитектура,
\textbf{Сравн. топ.} систематическое сравнение нескольких топологий
на едином стенде,
\textbf{Адапт. состава} динамическое изменение участников,
\textbf{Адапт. топ. runtime} переключение структуры графа в процессе
решения задачи,
\textbf{HITL} интеграция человека в контур управления как
полноценного участника графа,
\textbf{Безопасн.} учет устойчивости к скомпрометированному агенту
и каскадным сбоям. Знак \enquote{$+$} означает наличие свойства,
\enquote{$-$} --- его отсутствие, \enquote{$\pm$} --- частичную
реализацию.
\end{table}

\subsection{Что отсутствует в существующих работах}
\label{sec:rev:gap}

Из таблицы~\ref{tab:rev:methods} видно, что три столбца ---
систематическое сравнение топологий на едином стенде, адаптация
топологии во время решения задачи, явная роль человека в
структуре графа --- в существующих работах представлены лишь
эпизодически и в проанализированной выборке не встречаются
одновременно. Метод G-Designer~\cite{zhang2024}
систематически сравнивает топологии, но не адаптирует их во
время выполнения задачи. AutoGen~\cite{wu2023} позволяет подключать
человека, но не формализует его роль и не сравнивает разные
конфигурации. DyLAN~\cite{liu2023} адаптирует состав, но не
структуру связей. Ни один из методов не объединяет в одном
эксперименте сравнение нескольких топологий, механизм их
переключения во время решения и многомерное варьирование роли
человека.

Эта совокупность отсутствующих свойств и составляет
исследовательскую лакуну, заполнение которой является целью
настоящей работы. В частности, настоящая работа предлагает ---
во-первых, единый программный стенд, на котором сравниваются
пять различных топологий на едином корпусе из четырех классов
задач; во-вторых, мета-топологию, переключающую активную базовую
топологию в зависимости от фазы решения; в-третьих,
формализацию пяти ролей человека и экспериментальную оценку их
влияния на каждую топологию по полному факторному плану.

\subsection{Выводы по главе}
\label{sec:rev:summary}

Проведенный аналитический обзор позволяет сделать несколько
содержательных выводов. Мульти-агентные системы больших языковых
моделей являются активно развивающейся областью с устоявшимся
набором архитектурных образцов, однако коммуникационная топология
в них обычно фиксируется на этапе проектирования и не пересматривается
во время выполнения задачи. Существующие механизмы адаптации
ограничены либо составом команды, либо однократным офлайн-выбором
структуры графа, что недостаточно для задач со сложной фазовой
структурой. Безопасность мульти-агентных систем существенно зависит
от топологии, что подтверждает необходимость защитных правил при
ее динамическом изменении. Систематического исследования роли
человека в~графе агентов в~проанализированной выборке работ
не обнаружено. Среди смежных работ по~бенчмаркингу агентских систем
(AgentBench~\cite{agentbench2023}) и~по~маршрутизации между
языковыми моделями (RouteLLM~\cite{routellm2024},
FrugalGPT~\cite{frugalgpt2023}) формализованной адаптивной роли
человека также не введено.

На пересечении трех направлений --- сравнение топологий, их
динамическая адаптация, формализованная роль человека --- лежит
лакуна, заполнение которой ставится в качестве научной цели
настоящей работы. Эта цель формализуется в виде четырех
исследовательских вопросов в разделе~\ref{sec:methodology}.
